% ID: SEL-D0028-OB-E03
% ESTADO: cotejada. Segunda pasada visual completa y borrador compilado revisados.
% Fuente: Selectividad/matematicas_aplicadas/suplente junio/Examen-Matemáticas Aplicadas a las Ciencias Sociales.pdf
% SHA-256: a10146fce64a17e21450973c88d66d43571619781f76bc345b5a4b3dee67a16b
% Página: 2; opción B; ejercicio original 3.

Una caja contiene 3 bolas negras, 2 blancas y 1 roja. Se realiza el siguiente experimento aleatorio: ``Extraer de esa caja dos bolas al azar, una a continuación de otra sin reposición y anotar el color de las bolas en el orden en que han sido extraídas''.
\begin{enumerate}[label=\alph*)]
\item \textbf{(1 punto)} Describa el espacio muestral asociado a este experimento aleatorio.
\item \textbf{(1.5 puntos)} Indique la probabilidad de cada uno de los sucesos elementales del espacio muestral.
\end{enumerate}
